\chapter{Практический RLHF}
\label{ch:practical-rlhf}

\begin{quote}
\itshape
В теории нет разницы между теорией и практикой. На практике она есть. \upshape --- Йоги Берра (приписывается)
\end{quote}

\bigskip

В предыдущих главах была разработана теория моделей вознаграждения (Глава~\ref{ch:reward-models}), прямой оптимизации предпочтений (Глава~\ref{ch:dpo}) и группового относительного метода оптимизации стратегии (Глава~\ref{ch:grpo}). В каждой главе алгоритмы представлены в идеализированном виде: чистые функции потерь, элегантные выводы, гарантии сходимости при стандартных допущениях. Эта глава другая. Это инженерная глава\,---\,практическое руководство по всему, что находится между теоретическим алгоритмом и производственной системой RLHF.

Разрыв между теорией и практикой в RLHF огромен. Первопроходец, следующий правилу обновления PPO из Главы~\ref{ch:actor-critic}, обнаружит в течение нескольких часов, что процесс обучения расходится, модель вознаграждения эксплуатируется, память GPU исчерпана или модель коллапсировала в единственный повторяющийся режим. Каждый из этих режимов отказа был встречен\,---\,и решён\,---\,командами, стоящими за InstructGPT, Llama~2, Claude и DeepSeek, однако решения рассредоточены по блог-постам, техническим отчётам и племенным знаниям.

В этой главе эти знания собраны в единое связное изложение. Раздел~\ref{sec:rlhf-pipeline} описывает полный конвейер RLHF сквозным образом. Раздел~\ref{sec:collecting-preferences} охватывает искусство и науку сбора данных о предпочтениях. Раздел~\ref{sec:rm-practice} рассматривает обучение модели вознаграждения на практике. Раздел~\ref{sec:ppo-lm} подробно описывает реализацию PPO для языковых моделей, включая четырёхмодельную схему, таблицы гиперпараметров и соображения о смешанной точности. Раздел~\ref{sec:debugging-rlhf} является диагностическим руководством по наиболее распространённым режимам отказа. Раздел~\ref{sec:scaling-rlhf} рассматривает проблемы распределённых систем при масштабировании RLHF до больших моделей. Раздел~\ref{sec:safety-constraints} обсуждает ограничения безопасности и элайнмента, которые должны быть наложены поверх основного цикла обучения. Раздел~\ref{sec:case-studies} представляет краткие кейсы четырёх крупных работ по выравниванию.

% ===========================================================
\section{Конвейер RLHF сквозным образом}
\label{sec:rlhf-pipeline}

Прежде чем углубляться в детали каждого компонента, представим полный конвейер RLHF как единый поток. Конвейер состоит из шести этапов, каждый из которых подаёт данные в следующий.

\begin{center}
\begin{tikzpicture}[
  node distance=0.55cm and 0.3cm,
  stage/.style={
    rectangle, rounded corners=2pt, draw=black!70, fill=blue!6,
    minimum height=0.7cm, minimum width=1.6cm,
    font=\footnotesize, align=center, inner sep=3pt
  },
  arr/.style={-{Stealth[length=4pt]}, thick, black!60}
]
  \node[stage] (data)  {Данные\\предпочтений};
  \node[stage, right=of data]  (sft)   {SFT};
  \node[stage, right=of sft]   (rm)    {Модель\\вознаграждения};
  \node[stage, right=of rm]    (rl)    {RL\\(PPO/GRPO)};
  \node[stage, right=of rl]    (eval)  {Оценка};
  \node[stage, right=of eval]  (deploy){Развёртывание};

  \draw[arr] (data)  -- (sft);
  \draw[arr] (sft)   -- (rm);
  \draw[arr] (rm)    -- (rl);
  \draw[arr] (rl)    -- (eval);
  \draw[arr] (eval)  -- (deploy);

  % петля обратной связи
  \draw[arr, dashed] (eval.south) -- ++(0,-0.55) -| (data.south);
\end{tikzpicture}
\end{center}

\begin{description}[leftmargin=!,labelwidth=3.5cm]
  \item[Этап 1: Сбор данных] Собрать демонстрационные данные для SFT и попарные данные предпочтений для модели вознаграждения. Данные предпочтений могут поступать от экспертных аннотаторов, краудсорсеров или с помощью разметки с участием ИИ.
  \item[Этап 2: Supervised fine-tuning] Обучить базовую языковую модель на высококачественных демонстрациях для получения $\pi_{\mathrm{SFT}}$. Это отправная точка для всех последующих этапов и одновременно опорная стратегия $\pi_{\mathrm{ref}}$.
  \item[Этап 3: Обучение модели вознаграждения] Обучить скалярную модель вознаграждения $r_\psi$ на попарных предпочтениях, используя функцию потерь Bradley--Terry из Главы~\ref{ch:reward-models} (Уравнение~\eqref{eq:rm-loss}).
  \item[Этап 4: RL-обучение] Оптимизировать стратегию $\pi_\theta$ относительно $r_\psi$ с помощью PPO (Раздел~\ref{sec:ppo-lm}) или GRPO (Глава~\ref{ch:grpo}), со штрафом KL в сторону $\pi_{\mathrm{ref}}$ для предотвращения переоптимизации вознаграждения (Раздел~\ref{sec:overoptimization}).
  \item[Этап 5: Оценка] Оценить выровненную модель на отложенных бенчмарках предпочтений, тестах безопасности и человеческих оценках. Сравнить с $\pi_{\mathrm{SFT}}$ и предыдущими контрольными точками.
  \item[Этап 6: Развёртывание] Развернуть лучшую контрольную точку с мониторингом, логированием и фильтрами безопасности. Собирать новую обратную связь из продакшена для замыкания петли.
\end{description}

Пунктирная стрелка от оценки обратно к сбору данных представляет цикл \emph{итеративного выравнивания}: обратная связь из продакшена генерирует новые данные предпочтений, что запускает новый раунд обучения модели вознаграждения и стратегии. На практике крупные усилия по выравниванию запускают этот цикл несколько раз. InstructGPT (Ouyang et al., 2022) выполнил две полные итерации; Llama~2 (Touvron et al., 2023) --- пять.

\begin{remark}[Этапы не всегда последовательны]
\label{rem:pipeline-variations}
Не каждая система точно следует линейному шестиэтапному конвейеру. Некоторые команды обучают модель вознаграждения и модель SFT совместно (многозадачное обучение). Другие вообще пропускают отдельный этап SFT и начинают RL-обучение с предобученной базовой модели, компенсируя это мощной моделью вознаграждения. GRPO (Глава~\ref{ch:grpo}) устраняет модель вознаграждения для задач с верифицируемыми вознаграждениями. DPO (Глава~\ref{ch:dpo}) сворачивает этапы~3 и~4 в единый шаг обучения с учителем. Приведённый выше конвейер представляет наиболее распространённую архитектуру, однако практики должны воспринимать каждый этап как модуль, который можно заменить.
\end{remark}

% ===========================================================
\section{Сбор данных о предпочтениях}
\label{sec:collecting-preferences}

Качество данных о предпочтениях ограничивает качество всего последующего. Идеальный алгоритм RL, оптимизирующий шумную модель вознаграждения, обученную на небрежных аннотациях, даст посредственную стратегию. В этом разделе описаны передовые практики сбора данных о предпочтениях.

\subsection{Форматы аннотирования}
\index{preference data!annotation formats}

Два основных формата аннотирования предпочтений --- \emph{попарное сравнение} и \emph{оценка по шкале Лайкерта}.

\textbf{Попарное сравнение.}\quad По заданному промпту $x$ и двум ответам $y_A, y_B$ аннотатор выбирает лучший ответ. Это формат, принятый в модели Bradley--Terry (Определение~\ref{def:bt-model}). Его главное преимущество --- когнитивная простота: людям легче выносить относительные суждения, чем абсолютные. Основной недостаток --- неэффективность: сравнение $n$ ответов попарно требует $\binom{n}{2}$ аннотаций, тогда как оценки по Лайкерту требуют только $n$.

\textbf{Оценка по шкале Лайкерта.}\quad По заданному промпту $x$ и одному ответу $y$ аннотатор присваивает балл по фиксированной шкале (например, 1--7). Оценки по Лайкерту можно преобразовать в попарные предпочтения, считая ответ с более высокой оценкой предпочтительным, но это преобразование вносит шум, когда два рейтинга близки. Преимущество --- эффективность; недостаток --- индивидуальные рейтеры часто используют шкалу по-разному (``5'' одного аннотатора соответствует ``3'' другого).

На практике большинство крупных усилий по выравниванию используют попарное сравнение как основной формат. При сборе данных для Claude компания Anthropic использует четырёхбалльную шкалу (A намного лучше, A немного лучше, B немного лучше, B намного лучше), что позволяет включать поле в функцию потерь модели вознаграждения (Уравнение~\eqref{eq:rm-margin-loss}).

\subsection{Согласованность аннотаторов}
\index{inter-annotator agreement}

Когда несколько аннотаторов размечают одну и ту же пару, они нередко расходятся во мнениях. \emph{Межаннотаторское согласие} (IAA) измеряет согласованность человеческих суждений. Стандартная метрика --- каппа Коэна:\index{Cohen's kappa}
\[
\kappa = \frac{p_o - p_e}{1 - p_e},
\]
где $p_o$ --- наблюдаемая доля согласий, а $p_e$ --- ожидаемое случайное согласие. Для бинарных попарных предпочтений $p_e = 0.5$ и $\kappa = 2p_o - 1$. Значения $\kappa \geq 0.4$ обычно считаются умеренным согласием; $\kappa \geq 0.6$ --- существенным.

InstructGPT сообщал о IAA около 73\% ($\kappa \approx 0.46$) на своих данных предпочтений. Это согласуется с общим выводом о том, что аннотирование предпочтений --- по своей сути шумная задача\,---\,разные аннотаторы могут иметь подлинно различные предпочтения, а не просто случайные расхождения. Метод сглаживания меток cDPO (Уравнение~\eqref{eq:cdpo-loss}) явно учитывает этот шум.

\subsection{Качество данных против их количества}
\label{sec:data-quality-quantity}

Повторяющийся вопрос: стоит ли вкладываться в большее количество данных или в их лучшее качество. Эмпирические данные склоняются в пользу качества. Touvron et al.\ (2023) обнаружили, что качество модели вознаграждения Llama~2 достигало плато примерно после миллиона пар предпочтений, после чего наблюдались убывающие отдачи. Напротив, замена шумных краудсорсинговых меток экспертными аннотациями давала стабильное улучшение даже при одинаковом размере датасета.

\begin{remark}[Оценки стоимости данных предпочтений]
\label{rem:data-costs}
По состоянию на 2024 год типичные затраты на аннотирование попарных предпочтений составляют: краудсорсинг через такие платформы, как Scale AI или Surge AI, --- \$1--3 за сравнение; экспертные аннотации обученных лингвистов --- \$5--15 за сравнение; аннотации красной команды для безопасности --- \$10--25 за сравнение. Полный датасет для выравнивания из 100~000 пар предпочтений обходится примерно в \$100К--300К для краудсорсинговых данных и значительно дороже для экспертно-аннотированных. Эти затраты стимулируют растущий интерес к разметке с помощью ИИ, когда мощная языковая модель генерирует предварительные метки, которые затем проверяют люди-аннотаторы.
\end{remark}

\subsection{Распространённые ошибки}
\index{preference data!pitfalls}

\begin{itemize}[leftmargin=*]
  \item \textbf{Смещение по длине.} Аннотаторы склонны предпочитать более длинные ответы, даже когда дополнительная длина не добавляет ценности. Это переносится на модель вознаграждения, которая затем поощряет многословный вывод. Меры борьбы: инструктировать аннотаторов штрафовать излишнюю многословность; включать в набор для аннотирования пары с контролируемой длиной.
  \item \textbf{Позиционное смещение.} Когда два ответа представлены рядом, аннотаторы слегка предпочитают тот, что слева (или первый). Меры борьбы: рандомизировать порядок представления для каждой пары.
  \item \textbf{Неоднозначные промпты.} Когда промпт неясен, аннотаторы могут расходиться не потому, что имеют разные стандарты качества, а потому, что интерпретировали промпт по-разному. Меры борьбы: фильтровать промпты с высоким несогласием аннотаторов.
  \item \textbf{Утечка подхалимства.} Если аннотаторы предпочитают ответы, соглашающиеся с тезисами промпта (даже ложными), модель вознаграждения научится быть подхалимской. Меры борьбы: включать состязательные промпты с ложными тезисами и вознаграждать аннотаторов за выбор ответов, вежливо исправляющих пользователя.
\end{itemize}

% ===========================================================
\section{Обучение моделей вознаграждения на практике}
\label{sec:rm-practice}

В Главе~\ref{ch:reward-models} была выведена функция потерь Bradley--Terry и доказано, что модель вознаграждения сходится к истинной функции вознаграждения с точностью до константы. На практике ряд дополнительных соображений определяет, работает ли модель вознаграждения в действительности.

\subsection{Выбор архитектуры}
\index{reward model!architecture}

Модель вознаграждения $r_\psi(x, y)$ --- это, как правило, трансформер, принимающий на вход конкатенацию промпта и ответа и выдающий скалярный результат. Применяются две архитектурные стратегии.

\textbf{Та же архитектура, что и стратегия.}\quad Модель вознаграждения инициализируется из той же контрольной точки SFT, что и стратегия, при этом голова языкового моделирования заменяется линейной проекцией в скаляр. Этот подход используется InstructGPT и Llama~2. Преимущество --- модель вознаграждения начинает с сильным представлением языка; недостаток --- большой объём памяти, поскольку модель вознаграждения столь же велика, как и стратегия.

\textbf{Отдельная, меньшая модель.}\quad Модель вознаграждения использует меньший трансформер (например, 7B-модель вознаграждения для 70B-стратегии). Это снижает потребность в памяти, но рискует недостаточно подогнать распределение предпочтений. Gao et al.\ (2023) показали, что меньшие модели вознаграждения ведут к более быстрой переоптимизации (см. Пример~\ref{ex:overopt-scaling}), что говорит в пользу того, что ёмкость модели вознаграждения должна быть как можно больше в пределах бюджета.

\subsection{Устойчивость обучения и переобучение}
\label{sec:rm-overfitting}
\index{reward model!overfitting}

Модели вознаграждения склонны к переобучению, поскольку датасеты предпочтений относительно малы (от десятков тысяч до низких миллионов пар) по сравнению с корпусами предобучения (триллионы токенов). Стандартные меры противодействия включают:
\begin{itemize}[leftmargin=*]
  \item \textbf{Низкие скорости обучения.} Типичные значения --- от $1 \times 10^{-6}$ до $5 \times 10^{-6}$, что на порядок ниже скоростей обучения для SFT.
  \item \textbf{Ранняя остановка.} Мониторировать попарную точность на отложенном наборе и останавливаться, когда она начинает снижаться.
  \item \textbf{Сглаживание меток.} Заменить жёсткие метки Bradley--Terry на $1 - \epsilon$ и $\epsilon$ для некоторого малого $\epsilon$ (обычно $0.1$), как в формулировке cDPO (Уравнение~\eqref{eq:rm-label-smooth}).
  \item \textbf{Дропаут.} Применить дропаут ($p = 0.1$) к последним слоям модели вознаграждения.
\end{itemize}

\subsection{Калибровка и метрики оценки}
\label{sec:rm-eval}
\index{reward model!evaluation}

Модель вознаграждения может иметь высокую попарную точность и при этом давать плохо откалиброванные скалярные выходы. Ключевые метрики для мониторинга:
\begin{itemize}[leftmargin=*]
  \item \textbf{Попарная точность} (Уравнение~\eqref{eq:rm-accuracy}): доля отложенных пар предпочтений, где $r_\psi(x, y_w) > r_\psi(x, y_l)$. Значения 70--75\% типичны для выравнивания общего назначения; 80\%+ указывает на сильную модель вознаграждения.
  \item \textbf{Распределение вознаграждений.} Построить распределение баллов $r_\psi$ по большому набору ответов. Распределение должно быть унимодальным и примерно центрированным; бимодальное распределение говорит о том, что модель вознаграждения выучила ложный бинарный признак (например, длину ответа).
  \item \textbf{Согласие с человеческими ранжированиями.} На небольшом наборе промптов с четырьмя или более ответами, ранжированными экспертами, вычислить $\tau$ Кенделла между ранжированием модели вознаграждения и ранжированием экспертов.
\end{itemize}

\subsection{Когда переобучать}
\index{reward model!retraining}

Модель вознаграждения должна быть переобучена, когда стратегия достаточно далеко отдрейфовала от распределения, на котором была обучена модель вознаграждения. Практическое эмпирическое правило: мониторировать дивергенцию KL между текущей стратегией и опорной. Когда $\KL{\pi_\theta}{\pi_{\mathrm{ref}}}$ превышает пороговое значение (обычно 5--15 натов в зависимости от размера модели), собрать новые данные предпочтений от текущей стратегии и переобучить модель вознаграждения. Это итеративный цикл RLHF, описанный в Разделе~\ref{sec:rlhf-pipeline}.

% ===========================================================
\section{PPO для языковых моделей: детали реализации}
\label{sec:ppo-lm}

Применение PPO (Глава~\ref{ch:actor-critic}) к выравниванию языковых моделей --- не тривиальная задача. Структура MDP на уровне токенов при генерации языка (Определение~\ref{def:token-mdp}), огромное количество параметров и необходимость вычислять вероятности как под стратегией, так и под опорной моделью создают уникальный набор инженерных проблем. В этом разделе описана стандартная реализация.

\subsection{Четырёхмодельная схема}
\label{sec:four-models}
\index{RLHF!four-model setup}

Цикл обучения RLHF на основе PPO одновременно поддерживает четыре нейронные сети:
\begin{enumerate}
  \item \textbf{Модель стратегии} $\pi_\theta$: языковая модель, которая обучается. Генерирует ответы и вычисляет логарифмические вероятности.
  \item \textbf{Опорная модель} $\pi_{\mathrm{ref}}$: замороженная копия контрольной точки SFT. Используется для вычисления штрафа KL.
  \item \textbf{Модель вознаграждения} $r_\psi$: замороженная модель вознаграждения, обученная на данных предпочтений. Оценивает завершённые ответы.
  \item \textbf{Модель ценности} $V_\phi$: сеть-критик, оценивающая ожидаемое будущее вознаграждение с каждой позиции токена. Используется для оценки преимущества.
\end{enumerate}

При масштабе 7B параметров каждая модель требует приблизительно 14~ГБ в половинной точности (FP16/BF16). Хранение всех четырёх в памяти GPU одновременно требует не менее 56~ГБ\,---\,осуществимо на одной A100 (80~ГБ) при тщательном управлении памятью, но довольно плотно с учётом активаций и KV-кэша во время генерации. При масштабе 70B четыре модели требуют около 560~ГБ, что обусловливает необходимость многогрупповых GPU-установок (Раздел~\ref{sec:scaling-rlhf}).

\subsection{Реализация штрафа KL}
\label{sec:kl-implementation}
\index{KL penalty!implementation}

Штраф KL не позволяет стратегии слишком далеко отклоняться от опорной. Применяются две стратегии реализации.

\textbf{Вознаграждение KL на уровне токенов.}\quad На каждой позиции токена $t$ вычитаем штраф KL из вознаграждения:
\begin{equation}
\label{eq:per-token-kl}
\tilde{r}_t = \begin{cases}
  -\beta \log\dfrac{\pi_\theta(y_t \mid x, y_{<t})}{\pi_{\mathrm{ref}}(y_t \mid x, y_{<t})}, & t < T, \\[6pt]
  r_\psi(x, y) - \beta \log\dfrac{\pi_\theta(y_T \mid x, y_{<T})}{\pi_{\mathrm{ref}}(y_T \mid x, y_{<T})}, & t = T.
\end{cases}
\end{equation}
Терминальное вознаграждение при $t = T$ включает как оценку модели вознаграждения, так и штраф KL за последний токен. Эта формулировка распределяет сигнал KL по всем позициям токенов, что улучшает присвоение кредита. Она эквивалентна Уравнению~\eqref{eq:kl-token-reward} из Главы~\ref{ch:reward-models}.

\textbf{Штраф KL на уровне последовательности.}\quad Добавить суммарную дивергенцию KL уровня последовательности как единственный штраф в конце последовательности:
\[
\tilde{r} = r_\psi(x, y) - \beta \sum_{t=1}^{T} \log\frac{\pi_\theta(y_t \mid x, y_{<t})}{\pi_{\mathrm{ref}}(y_t \mid x, y_{<t})}.
\]
Это проще в реализации, но концентрирует весь сигнал KL на последнем токене, делая оценку преимущества менее информативной для промежуточных токенов.

На практике предпочтительна формулировка на уровне токенов.

\subsection{Оценка преимущества с помощью GAE}
\label{sec:gae-lm}
\index{GAE!language models}

Учитывая модифицированные вознаграждения на уровне токенов $\{\tilde{r}_t\}$ и оценки ценности $\{V_\phi(s_t)\}$, обобщённая оценка преимущества (GAE) из Главы~\ref{ch:actor-critic} равна
\begin{equation}
\label{eq:gae-lm}
\hat{A}_t = \sum_{l=0}^{T-t} (\gamma\lambda)^l \,\delta_{t+l}, \qquad \delta_t = \tilde{r}_t + \gamma V_\phi(s_{t+1}) - V_\phi(s_t),
\end{equation}
где $\gamma$ --- коэффициент дисконтирования, а $\lambda$ --- параметр GAE. Для RLHF языковых моделей стандартный выбор --- $\gamma = 1$ (без дисконтирования внутри ответа) и $\lambda \in [0.95, 1.0]$. Установка $\gamma = 1$ естественна, поскольку ``эпизод'' --- это один ответ: нет причин дисконтировать будущие токены внутри предложения.

\subsection{Формирование мини-батча и накопление градиентов}
\label{sec:minibatch}

Одна итерация PPO выполняется следующим образом:
\begin{enumerate}
  \item \textbf{Фаза сбора роллаутов.} Выбрать батч промптов из $\mathcal{D}$. Сгенерировать ответы с помощью $\pi_\theta$. Оценить каждый ответ с помощью $r_\psi$ и вычислить логарифмические вероятности на уровне токенов как под $\pi_\theta$, так и под $\pi_{\mathrm{ref}}$.
  \item \textbf{Вычисление преимущества.} Вычислить преимущества GAE с помощью модели ценности $V_\phi$.
  \item \textbf{Обновление PPO.} В течение $K$ эпох (обычно $K = 2$--$4$) перемешать данные роллаутов в мини-батчи и выполнить обновления градиента по отсечённой целевой функции PPO:
  \begin{equation}
  \label{eq:ppo-clip-lm}
  \mathcal{L}^{\mathrm{clip}}_t(\theta) = -\min\!\bigl(r_t(\theta)\,\hat{A}_t,\; \mathrm{clip}(r_t(\theta), 1{-}\epsilon, 1{+}\epsilon)\,\hat{A}_t\bigr),
  \end{equation}
  где $r_t(\theta) = \pi_\theta(y_t|x,y_{<t}) / \pi_{\theta_{\mathrm{old}}}(y_t|x,y_{<t})$ --- отношение важности, а $\epsilon$ --- параметр отсечения.
  \item \textbf{Обновление ценности.} Обновить $V_\phi$, минимизируя среднеквадратичную ошибку относительно целевых значений GAE-доходности.
\end{enumerate}

Когда памяти GPU недостаточно для полного мини-батча, используется \emph{накопление градиентов}: прямые и обратные проходы вычисляются на микро-батчах, а градиенты суммируются перед шагом оптимизатора. Это математически эквивалентно большему размеру батча, но требует только памяти для одного микро-батча.

\subsection{Обучение со смешанной точностью}
\index{mixed precision}

Современные реализации RLHF используют BF16 (bfloat16) для прямых проходов и FP32 для накопления градиентов и состояний оптимизатора. BF16 вдвое сокращает объём памяти по сравнению с FP32, сохраняя при этом динамический диапазон, необходимый для устойчивого обучения. Модель ценности иногда хранится в FP32 во избежание числовых проблем при вычислении преимущества, особенно когда вознаграждения охватывают широкий диапазон.

\subsection{Гиперпараметры}
\label{sec:ppo-hyperparams}

Таблица~\ref{tab:ppo-hyperparams} перечисляет типичные диапазоны гиперпараметров для RLHF на основе PPO для языковых моделей различных размеров.

\begin{table}[ht]
\centering
{\footnotesize

\resizebox{\textwidth}{!}{%
\begin{tabular}{@{}lll@{}}
\toprule
\textbf{Гиперпараметр} & \textbf{Типичный диапазон} & \textbf{Примечания} \\
\midrule
Скорость обучения (стратегия) & $5 \times 10^{-7}$ -- $2 \times 10^{-6}$ & Косинусный расписание, прогрев 3--10\% \\
Скорость обучения (ценность)  & $1 \times 10^{-6}$ -- $5 \times 10^{-6}$ & Часто в 2--5$\times$ больше LR стратегии \\
Коэффициент KL $\beta$ & $0.01$ -- $0.2$ & Начинать с малого, увеличивать при взрыве KL \\
Параметр отсечения $\epsilon$ & $0.1$ -- $0.2$ & $0.2$ наиболее распространён \\
GAE $\lambda$ & $0.95$ -- $1.0$ & $1.0$ типично для LM \\
Дисконт $\gamma$ & $1.0$ & Без дисконтирования внутри ответа \\
Эпохи PPO на батч $K$ & $2$ -- $4$ & Больше эпох рискует переобучением \\
Размер батча (промпты) & $64$ -- $512$ & Больше --- стабильнее \\
Размер мини-батча & $8$ -- $64$ & Ограничен памятью GPU \\
Максимальная длина ответа & $512$ -- $2048$ токенов & Зависит от задачи \\
\bottomrule
\end{tabular}%
}

}
\caption{Типичные гиперпараметры PPO для RLHF на языковых моделях.}
\label{tab:ppo-hyperparams}
\end{table}

% ===========================================================
\section{Отладка RLHF}
\label{sec:debugging-rlhf}

RLHF печально известна сложностью отладки. Тренировочные прогоны дороги, режимы отказа едва заметны, а взаимодействие четырёх нейронных сетей создаёт большое пространство возможных проблем. В этом разделе каталогизированы наиболее распространённые сбои и их диагностика.

\subsection{Взлом вознаграждения}
\index{reward hacking!symptoms}

\emph{Симптом:} Оценка модели вознаграждения быстро растёт, но качество генерируемого текста\,---\,по мнению людей\,---\,деградирует. Типичные проявления включают чрезмерно многословные ответы, повторяющиеся похвалы или формулировки-оговорки, и ответы, эксплуатирующие причуды форматирования.

\emph{Диагностика:} Строить в процессе обучения как прокси-вознаграждение (оценка модели вознаграждения), так и золотую метрику (например, процент побед над базовой линией по мнению людей). Взлом вознаграждения порождает расхождение: прокси-вознаграждение растёт, а золотая метрика останавливается или падает.

\emph{Меры устранения:}
\begin{itemize}[leftmargin=*]
  \item Увеличить коэффициент KL $\beta$. Это первый рубеж обороны.
  \item Переобучить модель вознаграждения на данных, сгенерированных текущей стратегией.
  \item Ансамблировать несколько моделей вознаграждения и использовать минимальную или медианную оценку.
  \item Добавить штраф за длину к вознаграждению для противодействия взлому путём многословия.
\end{itemize}

\subsection{Взрыв дивергенции KL}
\index{KL divergence!explosion}

\emph{Симптом:} Дивергенция KL $\KL{\pi_\theta}{\pi_{\mathrm{ref}}}$ быстро растёт (например, превышает 20 натов за несколько сотен шагов), и обучение становится нестабильным.

\emph{Диагностика:} Мониторировать KL на уровне токенов на каждой позиции. Если несколько позиций токенов имеют чрезвычайно высокий KL, возможно, стратегия научилась всегда производить определённый токен на этих позициях.

\emph{Меры устранения:}
\begin{itemize}[leftmargin=*]
  \item Снизить скорость обучения.
  \item Увеличить $\beta$ или перейти на адаптивный штраф KL, который динамически корректирует $\beta$ для поддержания целевого KL.
  \item Уменьшить количество эпох PPO на батч ($K$).
\end{itemize}

\subsection{Коллапс моды}
\index{mode collapse}

\emph{Симптом:} Стратегия генерирует почти одинаковые ответы на разные промпты. Метрики разнообразия (например, distinct $n$-gram, self-BLEU) резко падают.

\emph{Диагностика:} Выборочно сгенерировать 100 ответов стратегией на разнообразные промпты и вычислить попарные баллы BLEU. Средний попарный BLEU выше 0.5 --- сильный индикатор коллапса.

\emph{Меры устранения:}
\begin{itemize}[leftmargin=*]
  \item Увеличить $\beta$, чтобы стратегия оставалась ближе к разнообразной опорной.
  \item Увеличить температуру сэмплирования во время роллаутов.
  \item Добавить бонус энтропии к целевой функции PPO.
\end{itemize}

\subsection{Нестабильность обучения}
\index{RLHF!training instability}

\emph{Симптом:} Потеря стратегии дико колеблется, или оценка вознаграждения немонотонна с большими скачками между батчами.

\emph{Меры устранения:}
\begin{itemize}[leftmargin=*]
  \item Снизить скорость обучения и увеличить размер батча.
  \item Убедиться, что оценки преимущества нормализованы (нулевое среднее, единичная дисперсия) внутри каждого мини-батча.
  \item Проверить, не отстаёт ли модель ценности от стратегии; при необходимости обновлять её чаще.
  \item Проверить на числовые проблемы в BF16: переключить голову модели ценности на FP32.
\end{itemize}

\subsection{Диагностические метрики для мониторинга}
\label{sec:diagnostics}

Таблица~\ref{tab:diagnostics} суммирует ключевые метрики для логирования во время обучения RLHF.

\begin{table}[ht]
\centering
{\footnotesize

\resizebox{\textwidth}{!}{%
\begin{tabular}{@{}lll@{}}
\toprule
\textbf{Метрика} & \textbf{Здоровый диапазон} & \textbf{Тревожный сигнал} \\
\midrule
Среднее вознаграждение & Медленно растёт & Резкий скачок (взлом) \\
Дивергенция KL & $1$--$15$ натов & $> 20$ натов \\
Энтропия стратегии & Стабильна или медленно убывает & Внезапное падение (коллапс) \\
Доля отсечений & $0.05$--$0.15$ & $> 0.3$ (обновления слишком велики) \\
Потеря ценности & Убывает & Растёт (сбой модели ценности) \\
Средняя длина ответа & Стабильна & Монотонный рост (взлом многословием) \\
Прибл.\ норма градиента стратегии & Стабильна & Скачки (нестабильность) \\
\bottomrule
\end{tabular}%
}

}
\caption{Диагностические метрики для обучения RLHF. «Тревожный сигнал» означает значение, требующее расследования.}
\label{tab:diagnostics}
\end{table}

% ===========================================================
\section{Масштабирование RLHF}
\label{sec:scaling-rlhf}

Обучение RLHF при масштабе 70B+ параметров требует распределения вычислений по множеству GPU. Главная проблема состоит в том, что цикл обучения RLHF имеет две различные фазы\,---\,генерацию и обучение\,---\,с очень разными вычислительными профилями.

\subsection{Генерация против обучения}
\index{RLHF!generation phase}

\emph{Фаза генерации} (сбор роллаутов) является авторегрессионной: каждый токен зависит от всех предыдущих, что ограничивает параллелизм внутри одной последовательности. Основные оси параллелизма --- по последовательностям (параллелизм данных) и по слоям трансформера (тензорный параллелизм). Генерация ограничена памятью из-за KV-кэша, который растёт линейно с длиной последовательности.

\emph{Фаза обучения} (обновления PPO) является стандартным прямым-обратным проходом по последовательностям фиксированной длины. Она ограничена вычислениями и хорошо параллелизуется по батч-измерению (параллелизм данных) и по измерению модели (тензорный параллелизм или конвейерный параллелизм).

\subsection{Стратегии параллелизма}
\index{tensor parallelism}
\index{pipeline parallelism}

\begin{description}[leftmargin=!,labelwidth=3.5cm]
  \item[Параллелизм данных] Реплицировать модель на каждый GPU; каждый GPU обрабатывает разное подмножество батча. Градиенты синхронизируются с помощью AllReduce. Работает как для генерации, так и для обучения, но требует хранения полной копии модели на каждом GPU.
  \item[Тензорный параллелизм] Разбить отдельные матрицы весов по GPU внутри узла. Снижает потребление памяти на GPU, но требует высокоскоростных межсоединений (NVLink). Необходим для моделей, не умещающихся на одном GPU.
  \item[Конвейерный параллелизм] Разделить модель на стадии (группы слоёв), каждая из которых назначается отдельному GPU. Полезен для обучения в сочетании с планированием микро-батчей (например, 1F1B) для сокращения времени простоя.
  \item[Асинхронная генерация роллаутов] Разъединить фазы генерации и обучения: пул GPU генерирует роллауты асинхронно, пока другой пул выполняет обновления PPO на завершённых роллаутах. Этот подход используется в инфраструктуре RLHF OpenAI и в открытых фреймворках, таких как DeepSpeed-Chat и OpenRLHF.
\end{description}

\subsection{Оценки вычислительного бюджета}
\label{sec:compute-budgets}

Таблица~\ref{tab:compute-budget} содержит приблизительные оценки вычислений для одного раунда обучения RLHF при различных масштабах моделей, предполагая 50~000 промптов и 4 эпохи PPO.

\begin{table}[ht]
\centering
{\footnotesize
\begin{tabular}{@{}lccc@{}}
\toprule
\textbf{Размер модели} & \textbf{GPU (A100 80~ГБ)} & \textbf{Время (прибл.)} & \textbf{Стоимость (прибл.)} \\
\midrule
7B   & 8   & 12--24 часа & \$500--1~000 \\
13B  & 16  & 1--2 дня    & \$2~000--4~000 \\
70B  & 64  & 3--5 дней   & \$20~000--40~000 \\
\bottomrule
\end{tabular}
}
\caption{Приблизительный вычислительный бюджет для одного раунда обучения RLHF. Затраты основаны на облачной цене GPU \$2--3/GPU-час.}
\label{tab:compute-budget}
\end{table}

Эти оценки весьма вариативны и зависят от длины ответа, размера батча, числа эпох PPO и эффективности фреймворка распределённого обучения. Тем не менее они иллюстрируют ключевой момент: обучение RLHF дорого, но составляет малую долю стоимости предобучения. Типичный прогон предобучения для 70B-модели потребляет миллионы GPU-часов; один раунд RLHF --- тысячи.

% ===========================================================
\section{Ограничения безопасности и элайнмента}
\label{sec:safety-constraints}

Максимизация вознаграждения необходима, но недостаточна для безопасного развёртывания. В этом разделе обсуждаются дополнительные ограничения, которые должны быть наложены поверх основного цикла обучения RLHF.

\subsection{Красная команда}
\index{red-teaming}

\emph{Красная команда} --- это систематический процесс зондирования модели на предмет вредных выводов. Красная команда генерирует состязательные промпты, призванные вызвать токсичные, предвзятые или опасные ответы. Полученные сбои используются для формирования целевых данных предпочтений (где вредный ответ помечается как отвергнутый), которые затем включаются в следующий раунд обучения модели вознаграждения и стратегии.

Работа красной команды наиболее эффективна, когда она \emph{итеративна}: красная команда атакует последний контрольный пункт, сбои устраняются, и красная команда атакует снова. Ganguli et al.\ (2022) обнаружили, что работа красной команды с языковой моделью в течение 3--4 раундов снизила долю успешных атак на 75\% по сравнению с одним раундом.

\subsection{Интеграция Constitutional AI}
\index{Constitutional AI}

Constitutional AI (Bai et al., 2022) дополняет RLHF набором явных принципов («конституцией»), которым должна следовать модель. В стандартной реализации языковая модель критикует собственные выводы относительно конституции и генерирует исправленные ответы. Эти исправленные ответы используются как предпочтительный ответ в паре предпочтений, создавая синтетические данные предпочтений без разметки людьми. Конституционный подход может быть интегрирован в конвейер RLHF в качестве дополнительного источника данных для обучения модели вознаграждения.

\subsection{Полезность против безвредности}
\index{helpfulness-harmlessness tradeoff}

Устойчивое противоречие в выравнивании --- компромисс между полезностью и безвредностью. Модель, агрессивно обученная на безвредности, может отказываться отвечать на безобидные вопросы, которые поверхностно напоминают вредные (гиперотказ). Напротив, модель, обученная главным образом на полезности, может выполнять вредные запросы.

Стандартный способ снижения противоречия --- обучать отдельные модели вознаграждения для полезности и безвредности (или одну модель вознаграждения с отдельными головками) и комбинировать их во время RL-обучения:
\[
r(x, y) = r_{\text{helpful}}(x, y) + \alpha\,r_{\text{harmless}}(x, y),
\]
где $\alpha$ контролирует вес, придаваемый безвредности. Llama~2 использовала вариант этого подхода, при этом $\alpha$ возрастала в ходе обучения для акцентирования безопасности на более поздних стадиях.

\subsection{Когда останавливать обучение}
\label{sec:when-stop}

Не существует единого критерия для остановки обучения RLHF. На практике команды используют комбинацию:
\begin{itemize}[leftmargin=*]
  \item \textbf{Достижение плато в человеческой оценке.} Когда последовательные контрольные точки не показывают улучшения в человеческих оценках (например, рейтинги Эло в прямых сравнениях), дальнейшее обучение вряд ли поможет.
  \item \textbf{Исчерпание бюджета KL.} Некоторые команды устанавливают максимальный бюджет KL (например, 10 натов) и останавливают обучение при его достижении, полагая, что дальнейшее отклонение от опорной модели увеличивает риск переоптимизации.
  \item \textbf{Регрессия безопасности.} Если бенчмарки безопасности деградируют относительно предыдущей контрольной точки, обучение следует откатить.
\end{itemize}

% ===========================================================
\section{Кейсы}
\label{sec:case-studies}

В этом разделе кратко суммированы подходы к выравниванию четырёх влиятельных систем. Каждая представляет отдельную точку в пространстве проектных решений.

\subsection{InstructGPT (Ouyang et al., 2022)}
\index{InstructGPT}

InstructGPT стал первой крупномасштабной демонстрацией RLHF для языковых моделей. Конвейер следовал трёхэтапному рецепту: SFT на 13~000 демонстрациях, обучение модели вознаграждения на 33~000 пар предпочтений и дообучение PPO. Ключевые инженерные решения включали: использование 6B-модели вознаграждения для 175B-стратегии (экономия памяти в ущерб ёмкости модели вознаграждения); выполнение 2 полных итераций конвейера; использование коэффициента KL $\beta = 0.02$. Итоговая модель была предпочтительнее, чем 175B-базовая линия SFT, по мнению человеческих оценщиков в 71\% случаев.

\subsection{Llama 2 (Touvron et al., 2023)}
\index{Llama 2}

Усилия по выравниванию Llama~2 включали несколько инноваций. Команда собрала 1,4 миллиона пар предпочтений за 5 итераций цикла RLHF, с новыми данными от последней стратегии на каждой итерации. Они обучили две отдельные модели вознаграждения\,---\,одну для полезности, другую для безопасности\,---\,и объединили их с коэффициентом безопасности, который возрастал в ходе обучения. Примечательной техникой был \emph{rejection sampling} (отбор по отклонению): на каждой итерации команда генерировала $K$ ответов на промпт и выбирала лучший по модели вознаграждения, создавая синтетические SFT-данные для инициализации стратегии следующей итерации. Эта комбинация «RLHF + rejection sampling» оказалась эффективнее каждой техники по отдельности.

\subsection{Подход Anthropic (Bai et al., 2022)}
\index{Anthropic}

Работа Anthropic по выравниванию для Claude ввела Constitutional AI и сделала акцент на измерении безвредности. Ключевые вклады: механизм конституционной самокритики, генерирующий данные предпочтений из собственных ревизий модели; обучение с комбинированным вознаграждением полезность-безвредность; и обширная работа красной команды с итеративным устранением. Anthropic также стала пионером в использовании онлайн-RLHF, где модель вознаграждения периодически переобучается на предпочтениях, собранных от последней стратегии, что решает проблему сдвига распределения, от которой страдают офлайн-методы, такие как DPO (Раздел~\ref{sec:dpo-limitations}).

\subsection{Подход DeepSeek с GRPO (DeepSeek-AI, 2024)}
\index{DeepSeek}
\index{GRPO!practical use}

DeepSeek избрала другой путь, заменив PPO на Group Relative Policy Optimisation (GRPO, Глава~\ref{ch:grpo}). Для задач с верифицируемыми вознаграждениями\,---\,математики, программирования и фактических вопросов\,---\,GRPO полностью устраняет модель вознаграждения, используя сигналы корректности (например, частоту прохождения юнит-тестов, символьную верификацию) в качестве вознаграждения. Для открытых задач DeepSeek объединила GRPO с моделью вознаграждения. Ключевое инженерное преимущество --- устранение модели ценности: GRPO нормализует вознаграждения в группе ответов на один промпт, что даёт базовую линию без необходимости в отдельном критике. Это сокращает четырёхмодельную схему до двух моделей (стратегия и опорная), вдвое снижая потребность в памяти и упрощая инфраструктуру распределённого обучения.

% ===========================================================
\section{Резюме}
\label{sec:ch14-summary}

В этой главе было представлено практическое руководство по реализации RLHF в масштабе. Основные выводы:

\begin{enumerate}
  \item Конвейер RLHF состоит из шести этапов (сбор данных, SFT, обучение модели вознаграждения, RL-обучение, оценка, развёртывание), как правило, выполняемых итеративно.
  \item Качество данных о предпочтениях важнее их количества. Тщательные инструкции по аннотированию, рандомизированный порядок представления и внимание к смещениям (по длине, позиции, подхалимству) обязательны.
  \item Модели вознаграждения должны быть настолько большими, насколько позволяет вычислительный бюджет, обучаться с низкими скоростями обучения и мониторироваться на переобучение. Их необходимо переобучать, когда стратегия значительно дрейфует.
  \item Реализация PPO для языковых моделей требует одновременного хранения четырёх моделей в памяти. Штрафы KL на уровне токенов, GAE с $\gamma = 1$ и тщательное управление смешанной точностью являются стандартными инженерными решениями.
  \item Сбои RLHF\,---\,взлом вознаграждения, взрыв KL, коллапс моды\,---\,распространены и имеют хорошо известные диагностические паттерны и способы устранения. Мониторинг метрик из Таблицы~\ref{tab:diagnostics} обязателен.
  \item Масштабирование до больших моделей требует сочетания параллелизма данных, тензорного и конвейерного с асинхронной генерацией роллаутов.
  \item Ограничения безопасности\,---\,красная команда, Constitutional AI и компромисс полезность-безвредность\,---\,должны быть интегрированы в цикл обучения, а не добавлены после.
\end{enumerate}

Кейсы иллюстрируют, что не существует единственного «правильного» способа реализовать RLHF. InstructGPT использовал классический конвейер PPO. Llama~2 добавил rejection sampling и итеративный сбор данных. Anthropic ввела конституционную самокритику. DeepSeek заменил PPO на GRPO для верифицируемых задач. Каждый подход отражает различные компромиссы между инженерной сложностью, требованиями к данным и качеством выравнивания. Общая нить: успешное выравнивание требует не только правильного алгоритма, но и тщательного внимания к качеству данных, устойчивости обучения и оценке безопасности.

% ===========================================================
\section*{Упражнения}
\addcontentsline{toc}{section}{Упражнения}

\begin{exercise}
\label{ex:kl-budget}
Команда обучает 7B-языковую модель с RLHF на основе PPO. После 500 шагов обучения дивергенция KL $\KL{\pi_\theta}{\pi_{\mathrm{ref}}}$ достигла 18 натов, а средняя оценка модели вознаграждения выросла с 2.1 до 4.8.
\begin{enumerate}
  \item Находится ли дивергенция KL в здоровом диапазоне? Обоснуйте ответ, используя Таблицу~\ref{tab:diagnostics}.
  \item Команда замечает, что средняя длина ответа выросла с 120 токенов до 350 токенов. На какой режим отказа это указывает? Предложите два конкретных способа устранения.
  \item Предложите адаптивную схему для коэффициента KL $\beta$, которая поддерживает целевой KL $d_{\mathrm{target}} = 6$ натов. Запишите правило обновления.
\end{enumerate}
\end{exercise}

\begin{exercise}
\label{ex:four-model-memory}
Рассмотрим четырёхмодельную схему PPO (Раздел~\ref{sec:four-models}) для модели с 13B параметров.
\begin{enumerate}
  \item Оцените общую память GPU, необходимую для хранения всех четырёх моделей в BF16, без учёта активаций и KV-кэша.
  \item Команда хочет обучать на 8$\times$ A100 (по 80~ГБ). Достаточно ли параллелизма данных, или требуется тензорный параллелизм? Обоснуйте ответ.
  \item GRPO устраняет модель ценности. Сколько памяти это экономит и как меняются требования к параллелизму?
\end{enumerate}
\end{exercise}

\begin{exercise}
\label{ex:reward-hacking-detection}
Разработайте эксперимент для обнаружения взлома вознаграждения. Дана модель вознаграждения $r_\psi$ и доступ к пулу из 500 человеческих аннотаторов:
\begin{enumerate}
  \item Опишите протокол составления «золотого» оценочного набора, которого модель вознаграждения никогда не видела.
  \item Определите количественную метрику, позволяющую отличить подлинное улучшение от взлома вознаграждения. Точно опишите, что вы бы строили на графике и какой паттерн указывает на взлом.
  \item Команда замечает, что модель вознаграждения даёт оценку $+5.2$ ответу «С удовольствием помогу! Вот подробный, детальный и исчерпывающий ответ на ваш вопрос\ldots» вне зависимости от того, что следует дальше. Что это за тип взлома вознаграждения и как вы бы исправили модель вознаграждения?
\end{enumerate}
\end{exercise}

\begin{exercise}
\label{ex:annotation-design}
Вы разрабатываете конвейер аннотирования для обучения модели вознаграждения для чат-бота поддержки клиентов.
\begin{enumerate}
  \item Напишите инструкции по аннотированию (3--5 предложений) для попарного сравнения, учитывающие смещение по длине и подхалимство.
  \item Оцените стоимость сбора 50~000 пар предпочтений с помощью экспертных аннотаторов по \$10 за сравнение. Если у вас бюджет \$200~000, какие альтернативные стратегии сбора данных вы могли бы использовать для дополнения экспертно-аннотированных данных?
  \item Как вы измерите межаннотаторское согласие на вашем датасете? Каков минимальный уровень согласия, который вы требуете перед использованием данных для обучения модели вознаграждения?
\end{enumerate}
\end{exercise}

\begin{exercise}
\label{ex:gae-computation}
Рассмотрим ответ длиной $T = 4$ токена с вознаграждениями на уровне токенов $\tilde{r}_1 = 0.1$, $\tilde{r}_2 = -0.2$, $\tilde{r}_3 = 0.0$, $\tilde{r}_4 = 1.5$ и оценками ценности $V_1 = 0.8$, $V_2 = 0.6$, $V_3 = 0.9$, $V_4 = 1.0$, $V_5 = 0$ (терминальное состояние).
\begin{enumerate}
  \item Вычислите TD-невязки $\delta_1, \ldots, \delta_4$ при $\gamma = 1$.
  \item Вычислите преимущества GAE $\hat{A}_1, \ldots, \hat{A}_4$ при $\lambda = 0.95$.
  \item Повторите часть (б) при $\lambda = 1.0$ и сравните. Какая настройка даёт оценки преимущества с большей дисперсией и почему?
\end{enumerate}
\end{exercise}

\begin{exercise}
\label{ex:scaling-design}
Исследовательская лаборатория хочет запустить RLHF на 70B-параметрической модели с помощью PPO. У неё есть кластер из 64 GPU A100 (80~ГБ каждый), соединённых NVLink 400~Гбит/с внутри узлов из 8 GPU и InfiniBand 100~Гбит/с между узлами.
\begin{enumerate}
  \item Предложите стратегию параллелизма для фазы генерации и для фазы обучения. Укажите степень тензорного параллелизма, параллелизма данных и конвейерного параллелизма для каждой фазы.
  \item Фаза генерации является авторегрессионной и ограниченной памятью. Объясните, почему тензорный параллелизм помогает, несмотря на накладные расходы на коммуникацию.
  \item Команда рассматривает асинхронную генерацию роллаутов. Опишите потенциальную проблему устаревания: стратегия, использованная для генерации, может отличаться от той, что обновляется. При каких условиях это устаревание приемлемо?
\end{enumerate}
\end{exercise}

\begin{exercise}
\label{ex:case-study-comparison}
Сравните подходы InstructGPT, Llama~2 и DeepSeek к выравниванию по следующим измерениям:
\begin{enumerate}
  \item Количество моделей, требуемых в памяти GPU во время RL-обучения.
  \item Тип и количество использованных человеческих данных предпочтений.
  \item Стратегия для решения проблемы устаревания модели вознаграждения (сдвига распределения).
  \item Пригодность для задач с верифицируемыми вознаграждениями (например, математика).
\end{enumerate}
Представьте своё сравнение в таблице и напишите краткий абзац с рекомендацией, какой подход должен принять стартап с ограниченным вычислительным бюджетом для разработки ассистента общего назначения.
\end{exercise}
